\documentclass[11pt]{amsart}
\usepackage[a4paper,margin=22mm]{geometry}
\usepackage{llt-diagrams}
\usepackage[colorlinks=true]{hyperref}
\title{Diagram notation for restricted sums}
\author{}
\date{}
\begin{document}
\maketitle

Fix an ambient natural diagram, its strict edges, and a valuation $V(A)$
of a selected state. A cell diagram denotes the sum of $V(A)$ over all
states satisfying its restrictions. The diagonal lists the ordered vertex
labels; the cell in column $i$, row $j$ concerns the edge from the $i$th
label to the $j$th. Vertices suppressed from the picture retain their
specified exterior data. Unless we fix the exterior explicitly, we sum
over every permitted exterior selection as well.

\section*{Cell restrictions}
\begin{center}
\begin{tabular}{c@{\qquad}l}
\LLTDiagram{u,v}{} & natural edge, either selected or absent\\[8pt]
\LLTDiagram{u,v}{\LLTSelected{1}{2}} & selected normal edge\\[8pt]
\LLTDiagram{u,v}{\LLTAbsent{1}{2}} & normal edge, required absent from the selection\\[8pt]
\LLTDiagram{u,v}{\LLTForbidden{1}{2}} & no natural edge in the ambient diagram\\[8pt]
\LLTDiagram{u,v}{\LLTStrict{1}{2}} & strict edge, always selected\\[8pt]
\LLTDiagram{u,v}{\LLTSelected*{1}{2}} & selected normal edge required to be a cover\\[8pt]
\LLTDiagram{u,v}{\LLTStrict*{1}{2}} & strict edge required to be a cover
\end{tabular}
\end{center}

The underline imposes a cover condition in the \emph{full} selected relation,
including all exterior vertices. It does not merely exclude the displayed
two-step paths. A zero and a cross differ: zero restricts a selection,
whereas a cross removes the natural edge from the ambient diagram.
All unspecified cells represent optional natural edges. We draw a cross
where a pair is not natural.

\section*{Weights and extra conditions}
Gray cells carry no ascent weight. In the weighted convention, selected
normal edges contribute $q$ and their cells are white; strict edges
remain gray and contribute no $q$. Thus the two modes denote respectively
\[
 \sum_{A\text{ satisfying the picture}} V(A),\qquad
 \sum_{A\text{ satisfying the picture}}q^{|A\setminus S|}V(A).
\]
This weighted convention is the selected-state expansion after the LLT
substitution $q\mapsto q+1$. The exponent counts all selected normal edges,
including the exterior.
Color is redundant with the stated convention; entries still specify
the selection restrictions in a monochrome copy.
\[
 \LLTDiagram[name={\text{unweighted}},condition={|A|=s,\ \kappa(c)=\beta}]{a,b,c}
   {\LLTStrict{1}{3}}
 \qquad\qquad
 \LLTDiagram[weighted,name={\text{weighted}},condition={\kappa(c)=\beta}]{a,b,c}
   {\LLTStrict{1}{3}}
\]
For $V=E$ a fixed color refers to the canonical increasing HRV fibre chains.
Conditions below a diagram may also specify connectedness or component
membership. Such conditions are part of the sum, and are never inferred
from a small local picture.

\clearpage
\section*{The diamond}
For either $V=K$ or $V=E$, with any common fixed exterior and no cover
hypotheses, the following diagrams denote sums satisfying
\[
 \LLTSelectionTriangle{i,j,k}{1}{0}
 +\LLTSelectionTriangle{i,j,k}{0}{1}
 =\LLTSelectionTriangle{i,j,k}{0}{0}
 +\LLTSelectionTriangle{i,j,k}{1}{1}.
\]
The same identity holds with the color of a specified vertex fixed.
If all exterior choices are fixed as well, each picture denotes one
state contribution. The diamond is an equality of sums; it does not
assert that every two displayed states are joined by a flip.

\section*{NN: whole layers and weighted sums}
At an eligible NN contact, $x=w+1<z$. The exterior selection is summed
over its full permitted set. In the unweighted convention,
\[
 \LLTNNSource[condition={|A|=s}]{w,x,z}
 =\LLTNNTarget[condition={|A|=s}]{w,x,z}
 +\LLTNNTarget[condition={|A|=s-1}]{w,x,z}.
\]
Here $|A|$ includes strict edges. These are whole edge layers, each
partitioned into whole flip components. The identity does not assert
a recurrence for an arbitrary component without the correction argument.

Summing with normal-edge weights gives
\[
 \LLTNNSource[weighted]{w,x,z}
 =(1+q)\,\LLTNNTarget[weighted]{w,x,z}.
\]
Both sources and targets have the same number of strict edges, so the
two total-size layers give the two powers differing by one. The same
right-endpoint color can be fixed on every term.

\section*{The initial crossing slide}
The four displayed vertices satisfy $0<v<a<h$. The retained corner
covers are $0a$ and $vh$. The cell $va$ is a natural edge required to
be absent in these anchor endpoints.
\[
 \LLTDiagram[condition={\text{same exterior}}]{0,v,a,h}{%
   \LLTStrict*{1}{3}\LLTStrict*{2}{4}
   \LLTSelected{1}{2}\LLTAbsent{2}{3}\LLTAbsent{3}{4}
   \LLTForbidden{1}{4}}
 \quad\longleftrightarrow\quad
 \LLTDiagram[condition={\text{same exterior}}]{0,v,a,h}{%
   \LLTStrict*{1}{3}\LLTStrict*{2}{4}
   \LLTAbsent{1}{2}\LLTAbsent{2}{3}\LLTSelected{3}{4}
   \LLTForbidden{1}{4}}.
\]
This arrow denotes a path of actual flips between corresponding endpoint
states, with every exterior edge restored. It is not an equality of
their individual $K$ or $E$ values. Marking the selected maximal corners
strict makes their required status explicit.

\end{document}
